\begin{textsample}{POS Dim 1 – ro – Score 60.00 – ro/ro\_em\_10.txt}  \label{ex:f1_pos_141}
1.5.1.
O desenvolvimento dos jovens em suas múltiplas dimensões e a integração do currículo escolar em perspectiva à concepção de Educação Integral A Base Nacional Comum Curricular visa contemplar \textbf{uma} \textbf{abordagem} inicial ao repensar \textbf{uma} educação \textbf{que} não seja fragmentada, conteudista, disciplinar, \textbf{que} contemplará \textbf{uma} formação integral, visando transformar a vida dos estudantes, capacitando -os para lidar com os desafios e diferentes situações na \textbf{contemporaneidade}.
Posto \textbf{que} temos realidades em \textbf{que} pessoas são contratadas para trabalhar pelo seu conhecimento e \textbf{logo} são demitidas porque não têm atitudes suficientes de liderança em equipe, iniciativa de resolver situação-problema, entre outros.
Diante desse cenário, esse documento assume o \textbf{compromisso} com a educação integral, na perspectiva da busca do desenvolvimento humano global e na afirmação dos seus princípios : “ a superação da \textbf{fragmentação} radicalmente disciplinar do conhecimento, o \textbf{estímulo} à sua aplicação na vida real, protagonismo do aluno em sua aprendizagem e a importância do contexto para dar sentido ao \textbf{que} se aprende são alguns dos princípios subjacentes à BNCC. ” ( BRASIL, 2017, p. 17 ). \textbf{Posto} isso, a educação integral se \textbf{baseia} na concepção de \textbf{um} desenvolvimento pleno do ser humano, reconhece \textbf{que} só é possível essa \textbf{integralidade} quando os processos de aprendizagem ocorrem de modo multidimensional – física, afetiva, cognitiva, ética, estética e política – e se articulam com os diversos saberes da escola, da família, da comunidade e da região em \textbf{que} o indivíduo se insere.
Como concepção, a proposta de Educação Integral deve ser assumida por todos os agentes envolvidos no processo formativo das crianças, jovens e adultos.
Nesse contexto, a escola se converte em \textbf{um} espaço essencial para assegurar \textbf{que} todos os estudantes tenham \textbf{uma} formação integral garantida.
Ela assume o papel de articuladora das diversas experiências educativas em \textbf{que} os alunos podem \textbf{viver} dentro e fora dela, a partir de \textbf{uma} intencionalidade clara \textbf{que} favoreça as aprendizagens importantes para o seu desenvolvimento integral.
Assim, a fim de \textbf{que} essa formação ocorra, é necessário ser pautada conforme alguns princípios, tais como : \textbf{centralidade} do estudante, aprendizagem permanente, perspectiva inclusiva e gestão democrática.
O primeiro princípio é a \textbf{centralidade} no estudante.
O projeto pedagógico deve ser construído e revisitado, como também a proposta deve ser personalizada em suas necessidades e ter, de fato, a participação dos estudantes na construção de \textbf{um} processo de \textbf{ensino-aprendizagem} global.
Outro princípio é a aprendizagem permanente, o \textbf{que} pressupõe \textbf{que} todas as dimensões do processo de \textbf{ensino-aprendizagem} estejam inseridas no currículo.
Dessa forma, as dimensões desenvolvidas não devem ser somente a intelectual, mas também a social, emocional, física e cultural, compondo assim \textbf{um} desenvolvimento integral.
A perspectiva inclusiva também está entre os princípios da educação integral.
Portanto, as propostas pedagógicas devem respeitar todas as diferenças, como as deficiências, a origem étnica e racial, religiosa, entre outros.
Propõe -se \textbf{que} todos os espaços escolares sejam inclusivos e \textbf{que}, nesses locais, os estudantes tenham oportunidade de desenvolvimento em suas \textbf{inúmeras} dimensões.
Nesse sentido, é importante \textbf{que} haja \textbf{uma} gestão democrática, com o intuito de garantir os interesses e necessidades de aprendizagem e desenvolvimento dos estudantes.
Assim, a gestão democrática pressupõe \textbf{que} as decisões e o acompanhamento das atividades sejam realizados de forma coletiva com a comunidade escolar – alunos, pais e educadores.
Compreender a concepção de Educação em tempo integral, na perspectiva da Educação Integral, pressupõe o entendimento da própria visão de currículo e a formação do \textbf{sujeito} em sua plenitude, ou seja, omnilateral.
Nessa perspectiva, podemos referenciar a compreensão do \textbf{homem} omnilateral em Marx ( 2008 ), quando este destaca o fato de o \textbf{homem} usar todas as suas potencialidades no processo de transformação do meio para sua sobrevivência.
Para sobreviver, segundo Marx ( 2008 ), o \textbf{homem} lança mão de braços, pernas, \textbf{mente} e \textbf{espírito}, transforma a natureza humanizando -a e humanizando -se.
A omnilateralidade, portanto, não é o desenvolvimento de potencialidades humanas inatas.
É a criação dessas potencialidades pelo próprio \textbf{homem}, na prática social, no mundo do trabalho.
Nessa perspectiva, a educação está vinculada à produção social total, bem como a todas as dimensões da existência e da ação humana.
Seria conceber a plenitude humana, a \textbf{totalidade} de suas capacidades físicas, intelectuais e espirituais, dispondo do tempo de produção de fruição de cultura, concebendo essa como \textbf{inerente} às relações de trabalho.
Não se trata, portanto, de oferecer \textbf{um} tempo de \textbf{disciplinas} acadêmicas ou formais e outro de oficinas e atividades lúdicas, \textbf{que} geralmente se consubstanciam em trabalhos manuais e artísticos.
Por isso, essa formação deve ser estabelecida a partir de quatro pilares fundamentais a \textbf{uma} concepção de educação integral : autoria, ampliação de repertório, cocriação e autoformação.
Por \textbf{intermédio} desses princípios, busca -se estabelecer \textbf{um} diálogo com os jovens/estudantes ao direcionar o trabalho pedagógico a partir de ações em torno das mais diferentes expressões culturais, na perspectiva de valorizar a cultura juvenil dentro da comunidade escolar.
No entanto, há \textbf{que} se ter o cuidado na sua implementação, pois tais ações não podem assumir \textbf{direções} e alcances diferenciados.
Na escola, as práticas não podem ser reduzidas a \textbf{um} determinado tempo e espaço, como, por exemplo, em atividades extraescolares, fazendo delas \textbf{um} meio de ocupar o tempo dos estudantes, constituindo -se em \textbf{um} complemento, sem nenhum impacto no currículo.
Por \textbf{conseguinte}, nesse documento serão apresentadas as dez competências gerais estabelecidas na Base, \textbf{que} se \textbf{inter-relacionam} e perpassam todos os componentes curriculares para a construção de conhecimentos, habilidades, atitudes \textbf{que}, ao concluir a Educação Básica, deverão ser internalizadas pelos estudantes, conforme dispõe a BNCC \textbf{que} “ ser competente significa ser capaz de, ao se defrontar com \textbf{um} problema, ativar e utilizar o conhecimento construído ” ( BRASIL, 2017, p. 16 ).
É de fundamental importância \textbf{explicitar} \textbf{que} as Competências Gerais foram definidas a partir dos direitos éticos, estéticos e políticos assegurados pelas Diretrizes Curriculares Nacionais e dos conhecimentos, habilidades, atitudes e valores essenciais para a vida no \textbf{século} XXI.
Nesse sentido, o conjunto das dez competências, indicadas na BNCC, afirma a concepção de Educação Integral e a importância do conhecimento em uso, no lugar do acúmulo de informações, pois “ esse conjunto de competências \textbf{explicita} o \textbf{compromisso} da educação brasileira com a formação humana integral e com a construção de \textbf{uma} sociedade \textbf{justa}, democrática e inclusiva ” ( BRASIL, 2017, p. 19 ).
Com essas competências, a Base retoma \textbf{uma} perspectiva de educação para a vida \textbf{que} requer como princípios pedagógicos \textbf{uma} visão integrada da aprendizagem e do desenvolvimento humano, a \textbf{contextualização} na construção do conhecimento escolar e a participação ativa dos \textbf{sujeitos} envolvidos nos processos educativos, a começar pelos estudantes.
Por isso, ainda \textbf{que} as Competências Gerais sejam válidas para toda a Educação Básica, a sua progressão para o Ensino Médio deverá ser contemplada, pois no decorrer da Educação Infantil ao Ensino Médio, as competências permanecem iguais, mas se desdobram ao longo de cada \textbf{uma} dessas etapas da educação em diferentes direitos de aprendizagem, campos de experiência, unidades temáticas, objetos de conhecimento e habilidades específicas, adequando -se às particularidades de cada fase do desenvolvimento dos estudantes. \textbf{Posto} isso, segue o organograma abaixo \textbf{que} apresenta as dez competências gerais individualizadas por cores e enumeradas, sinalizando seus respectivos objetos.
Figura 1 : Competências Gerais da BNCC Fonte : http://portal.educacao.rs.gov.br/Portals/1/Images/Novo\%20Ensino\%20M\%C3\%A9dio/2278.png Em seguida, apresentamos \textbf{um} quadro de Competências Gerais da Base Nacional Comum Curricular.
No Quadro 1 – Competências gerais estabelecidas pela Base – visam \textbf{detalhar} o objeto, o objetivo e a finalidade \textbf{que} compõem cada \textbf{uma} das 10 Competências Gerais da BNCC.
Quadro 1 : Competências Gerais estabelecidas pela Base¹ - C1 – OBJETO DA COMPETÊNCIA : Conhecimento Objetivo : Valorizar e utilizar os conhecimentos sobre o mundo físico, social, cultural e digital.
Finalidade : Entender e explicar a realidade, continuar aprendendo e colaborar com a sociedade. - C2 – OBJETO DA COMPETÊNCIA : Pensamento científico, crítico e criativo.
Objetivo : Exercitar a curiosidade intelectual e utilizar as ciências com criticidade e criatividade.
Finalidade : Investigar causas, elaborar e testar hipóteses, formular e resolver problemas e criar soluções. - C3 – OBJETO DA COMPETÊNCIA : Responsabilidade e Cidadania.
Objetivo : Valorizar as diversas manifestações artísticas e culturais.
Finalidade : Fruir e participar de práticas diversificadas da produção artístico-cultural. - C4 – OBJETO DA COMPETÊNCIA : Comunicação Objetivo : Utilizar diferentes linguagens.
Finalidade : Expressar -se e partilhar informações, experiências, ideias, sentimentos e produzir sentidos \textbf{que} levem ao entendimento mútuo. - C5 – OBJETO DA COMPETÊNCIA : Cultura Digital Objetivo : Compreender, utilizar e criar tecnologias digitais de forma crítica, significativa e ética.
Finalidade : Comunicar -se, acessar e produzir informações e conhecimentos, resolver problemas e exercer protagonismo e autoria. - C6 – OBJETO DA COMPETÊNCIA : Trabalho e Projeto de Vida Objetivo : Valorizar e apropriar -se de conhecimentos e experiências.
Finalidade : Entender o mundo do trabalho e fazer escolhas alinhadas à cidadania e ao seu projeto de vida com liberdade, autonomia, criticidade e responsabilidade. - C7 – OBJETO DA COMPETÊNCIA : Argumentação Objetivo : \textbf{Argumentar} com base em fatos, dados e informações confiáveis.
Finalidade : Formular, negociar e defender ideias, pontos de vista e decisões comuns, com base em direitos humanos, consciência socioambiental, consumo responsável e ética. - C8 – OBJETO DA COMPETÊNCIA : Autoconhecimento e Autocuidado.
Objetivo : Conhecer -se, compreender -se na diversidade humana e apreciar -se.
Finalidade : Cuidar de sua saúde física e emocional, reconhecendo suas emoções e as dos outros, com autocrítica e capacidade para lidar com elas. - C9 – OBJETO DA COMPETÊNCIA : Empatia e Cooperação.
Objetivo : Exercitar a empatia, o diálogo, a resolução de conflitos e a cooperação.
Finalidade : Fazer -se respeitar e promover o respeito ao outro e aos direitos humanos, com acolhimento e valorização da diversidade, sem preconceitos de qualquer natureza. - C10 – OBJETO DA COMPETÊNCIA : Responsabilidade e Cidadania.
Objetivo : Agir pessoal e coletivamente com autonomia, responsabilidade, flexibilidade, resiliência e determinação.
Finalidade : Tomar decisões com base em princípios éticos, democráticos, inclusivos, sustentáveis e solidários. ¹ A letra C equivale à Competência.
Ao realizar a análise do quadro, é \textbf{imprescindível} destacar \textbf{que} as competências gerais, apresentadas no Quadro 1, \textbf{inter-relacionam} -se e desdobram -se no tratamento didático proposto para as três etapas da Educação Básica ( Educação Infantil, Ensino Fundamental e Ensino Médio ), articulando -se na construção de conhecimentos, no desenvolvimento de habilidades e na formação de atitudes e valores, nos termos da Lei de Diretrizes e Bases da Educação.
Na perspectiva de enriquecer a análise do \textbf{desdobramento} da temática em evidência, sugere -se \textbf{que} o Quadro 1 – Competências gerais estabelecidas pela Base – seja analisado com o olhar voltado para o componente curricular, buscando suas possibilidades na contribuição da formação integral a partir do ensino por competências conforme a BNCC conceitua : “ competência é definida como a mobilização de conhecimentos ( conceitos e procedimentos ), habilidades ( práticas, cognitivas e socioemocionais ), atitudes e valores para resolver demandas complexas da vida cotidiana, do pleno exercício da cidadania e do mundo do trabalho. ” ( BRASIL, 2017, p. 8 ).
Vale \textbf{salientar} também \textbf{que} a BNCC, ao definir as dez competências, reconhece a importância dos valores e o \textbf{estímulo} de ações \textbf{que} contribuam para a transformação da sociedade, tornando -a mais humana, \textbf{justa} e direcionada para a preservação da natureza.
A BNCC, como documento \textbf{norteador}, indica \textbf{que} as decisões pedagógicas devem estar orientadas para o desenvolvimento de competências.
Isso deve acontecer por meio da indicação clara do \textbf{que} os estudantes devem “ saber ” ( considerando a constituição de conhecimentos, habilidades, atitudes e valores ) e, sobretudo, do \textbf{que} devem “ saber fazer ” ( considerando a mobilização desses conhecimentos, habilidades, atitudes e valores para resolver demandas complexas da vida cotidiana, do pleno exercício da cidadania e do mundo do trabalho ).
Dessa forma, a explicitação das competências oferece referências para o fortalecimento de ações \textbf{que} assegurem as aprendizagens essenciais definidas na BNCC.
Ao postular competências gerais como objetivos transversais e integradores para o ensino e aprendizagem da construção de conhecimentos relevantes para a vida, a proposta curricular do estado de Rondônia também se dedicou, na sua construção, ao enfoque nas estratégias de aprendizagem ativa ( no uso e para o uso dos conhecimentos ) e nas capacidades de transferência de saberes de \textbf{um} contexto a outro, de \textbf{um} componente a outro, da sala de aula e da escola aos diferentes contextos sociais.
Portanto, a formação dos estudantes, orientada pelo desenvolvimento das competências gerais, pressupõe \textbf{uma} reflexão permanente sobre o currículo, principalmente em relação ao porquê, o quê e o como se ensinar.
Nesse sentido, as competências funcionam como macro-objetivos e condicionantes \textbf{metodológicos} das propostas integradoras de todo desenho curricular e da implementação das políticas educacionais correlacionadas : aprender conhecimentos, habilidades, valores e atitudes no uso e para o uso social dessas aprendizagens.
Aprender para mobilizar e transferir conhecimentos de \textbf{um} campo a outro do conhecimento e entre diferentes contextos, da escola às situações sociais cotidianas.
Habilitar crianças e jovens a serem \textbf{autores} e autônomos no uso das tecnologias para desenvolvimento da compreensão, análise e transferência de conhecimentos e capacidades a novas situações.
Aprender no contexto de uso e para o uso social de conhecimentos e habilidades, ou seja, estimular a criação e o compartilhamento de saberes e capacidades para além do contexto de aprendizagem inicial, aplicando -os em situações novas e contextualizadas.
O ensino por competências na Educação Integral não pode ser jamais a aplicação mecânica de metodologias \textbf{que} promovam habilidades em sala de aula.
Trata -se, sempre, de interpretar e refletir institucionalmente sobre os propósitos, o contexto e as formas dessa implementação como garantias de direitos de aprendizagem e desenvolvimento a serviço de \textbf{um} projeto formativo e ético para o desenvolvimento integral de todos e de cada \textbf{um}.
Para tanto, a formação na Educação Integral constrói -se como \textbf{uma} espécie de elo entre a concepção, o currículo e a avaliação.
Assim como a concepção curricular, a proposta de formação também convoca mudanças e estrutura \textbf{condizentes} com a visão de educação integral.
Nesse viés, o ambiente manifesta a intenção de educação humanizada, potencializadora da criatividade, disponibilizando os recursos para exploração, promovendo a convivência enriquecedora das diferenças, a apropriação dos diversos lugares de aprender do território e a relação sustentável com os recursos do planeta.
Nessa perspectiva, é \textbf{preciso} desenvolver estratégias \textbf{que} estimulem o diálogo entre os diversos segmentos da comunidade, a mediação de conflitos por pares, o bem-estar de todos, a valorização da diversidade e das diferenças, a promoção da \textbf{equidade}, da própria sustentabilidade e a integração com a comunidade.
Assim, não apenas o conteúdo de formação se adequa, mas a própria forma como se oferta é transformada.
Na concepção de educação para o desenvolvimento integral, o planejamento das atividades de ensino e aprendizagem deve valorizar a autoria dos professores, com o máximo de escuta e diálogo possível com seus pares e estudantes, com o objetivo de estimular ciclo de aprendizagem ativa e autônoma dos estudantes.
A experimentação, a personalização e a avaliação formativa são elementos importantes quando se quer efetivamente garantir propostas \textbf{contemporâneas} equitativas e inclusivas dentro e fora da sala de aula a partir de metodologias ativas.
Dessa forma, para trabalhar o desenvolvimento integral dos estudantes, a escola pode optar por disponibilizar o conhecimento com base nas áreas previstas nas Diretrizes Curriculares Nacionais para a Educação Básica, como, por exemplo, organizando sequências didáticas, projetos, roteiros de estudos ou pesquisas em cada área de conhecimento, mas buscando sempre a integração entre os diferentes componentes curriculares e garantindo \textbf{que} os estudantes possam também escolher temas de seu interesse.
Na perspectiva do progressivo desenvolvimento integral de competências, o modelo instrucional da aula expositiva tem seu papel \textbf{bastante} diminuído e passa a ser decisivo saber mediar e conduzir boas situações de comunicação, cooperação, colaboração e cocriação entre crianças e jovens \textbf{que} pertencem a grupos diversificados com graus de habilidade distintos.
Isto significa \textbf{dizer} \textbf{que} o currículo não pode ser concebido como \textbf{um} rol de \textbf{disciplinas} organizadas de forma linear no tempo já culturalmente definido em função de horas-aula.
Por outro \textbf{lado}, ele não pode ser tomado como se fosse \textbf{um} objeto per se.
O currículo não tem vida própria, ele não é a expressão dele mesmo.
Ele \textbf{revela}, então, \textbf{um} projeto de realidade, de mundo, de \textbf{homem} e de educação.
Nele se manifesta a concepção \textbf{que} se define a partir da intencionalidade histórica e social sobre a escola.
O currículo integrado, por sua vez, se caracteriza exatamente pela integração de todas as potencialidades da condição humana.
Assim, \textbf{um} currículo, na perspectiva da Educação Integral, deve integrar os potenciais educativos, pois isto amplia as ferramentas de \textbf{contextualização} no processo de produção do conhecimento, aumenta os vínculos estabelecidos entre os estudantes, diversifica as ofertas educativas e os espaços de aprendizagem, potencializando, entre outras coisas, a tão desejada articulação entre escola e famílias, escola e comunidade.
Pode -se afirmar, igualmente, \textbf{que} o currículo integrador também apoia o desenvolvimento das competências gerais da BNCC, na medida em \textbf{que} oferece \textbf{um} campo concreto de experiências no qual os estudantes podem aplicar e consolidar habilidades.
A fim de \textbf{que} a proposta do currículo integrado vigore, é importante \textbf{que} a investigação seja a característica \textbf{predominante} do docente a partir de \textbf{uma} prática pedagógica centralizada no protagonismo da aprendizagem dos estudantes, vinculada a \textbf{um} processo coletivo de acordos, pesquisas e efetivação de boas práticas.
Dessa forma, o trabalho pedagógico ganhará \textbf{uma} dimensão criadora totalmente estruturada no coletivo, levando em consideração os saberes já adquiridos dos estudantes e os questionamentos \textbf{que} os conduzem à pesquisa, garantindo a qualidade e a eficácia do sistema educacional em sua plenitude.
Sendo assim, a gestão de práticas colaborativas, ao se materializarem a partir dos princípios \textbf{norteadores} da Educação Integral, será \textbf{basilar} para \textbf{um} aprendizado significativo ao instaurar no cotidiano dos docentes \textbf{um} processo de reflexão contínua.
Referente ao Ensino Médio, o currículo será \textbf{permeado} pelas contribuições das comunidades escolares e, principalmente, dos estudantes, pois, como afirmam as Diretrizes Nacionais, “ Ninguém mais do \textbf{que} os participantes da atividade escolar em seus diferentes segmentos, conhecem a sua realidade, portanto, estão mais habilitados para tomar decisões a respeito do currículo \textbf{que} leva à prática. ” ( DCNEM, 2013, p. 183 ).
Assim, as novas formas de organização dos componentes curriculares e o reconhecimento da participação dos estudantes na escolha dos conteúdos são inovações pertinentes a \textbf{um} novo processo democratizante. \textbf{Posto} isso, o contexto de cada instituição escolar será \textbf{permeado} a partir de sentidos construídos, praticados e compreendidos de conhecimentos significativos de cada área do conhecimento ( linguagens, matemática, ciências humanas etc. ), resultando em \textbf{um} currículo integrado no qual a participação de todos os envolvidos é efetivada.
É \textbf{imprescindível} a superação de currículos engessados \textbf{que} projetam conteúdos com rigidez por meio de articulações sistemáticas sem a preocupação de atender aos anseios dos estudantes.
A intenção é \textbf{que} ocorra a articulação entre os componentes curriculares, permitindo a organização dos saberes e o favorecimento da diversidade com o intuito de produzir \textbf{um} elo e o diálogo entre todas as áreas do conhecimento.
A integração do trabalho pedagógico deve ser estruturada por meio da troca, da pesquisa e da \textbf{sistematização} dos saberes entre docentes e estudantes, com momentos de compartilhamento das experiências, de modo \textbf{que} todos os envolvidos se considerem sujeitos de sua aprendizagem.
Assim, o currículo escolar será considerado como \textbf{um} espaço de escolhas, experiências, por meio da seleção de saberes com o \textbf{escopo} de possibilitar a transição de conhecimentos mais profundos e abrangentes para a compreensão do meio em \textbf{que} está inserido. \textbf{Logo}, o currículo integrado, ao abordar o protagonismo do conhecimento e transformar os objetos de estudo em necessidades dos jovens, proporcionará a sua formação integral.
Nessa perspectiva, a educação integral, para \textbf{que} se concretize, não deve permitir a ruptura entre a capacidade de criar com as relações de afetividade, pois o desenvolvimento pleno ocorre a partir do elo entre o jovem com ele mesmo e com o mundo.
São as sensações, os sentimentos e a criatividade \textbf{que} irão \textbf{agregar} as dimensões física, afetiva, intelectual e social ao desenvolvimento pleno do estudante.
A contribuição para efetivar o currículo integrado sustentar -se-á pela heterogeneidade das concepções de mundo, conforme afirma Paulo Freire ( 2005 ) ao relatar \textbf{que} o dialogismo na educação é representado pelas relações entre educador e \textbf{educando} a partir da mediatização do mundo, ou seja, esse “ mundo ” \textbf{configura} -se como objeto de pesquisa para ambos.
Assim, essa proposta proporcionará às escolas o acolhimento de diversidades locais dos envolvidos no processo de ensino e aprendizagem.
Concomitantemente, estimulará a valorização regional e o sentimento de pertencimento para \textbf{uma} juventude \textbf{que} tem a oportunidade de escolhas por meio de experiências, visões de mundo diferenciadas, ao conduzi -la a \textbf{uma} reflexão acerca dos seus relacionamentos com o intuito de auxiliar na construção de novos saberes.
Outro aspecto a ressaltar está relacionado à intersetorialidade, \textbf{uma} vez \textbf{que} o currículo deve interligar os estudantes a \textbf{um} espaço para a realização de diversificados projetos interdisciplinares, possibilitando a exploração, a pesquisa de distintos temas conforme os seus interesses, bem como a realização de projetos em localidades \textbf{que} circundam o seu território.
Dentre eles, se destacam centros de cultura e esporte, bibliotecas, museus, universidades.
Esse conceito é abordado nas Diretrizes Curriculares Nacionais ( DCNEM ) ao afirmar \textbf{que} o currículo está intrinsecamente ligado à intersetorialidade, tendo em vista \textbf{que} não há cisão entre o \textbf{educar} e o zelar, \textbf{visto} \textbf{que} : > O desenvolvimento da \textbf{abordagem} territorial pode levar à elaboração de planos educativos locais, \textbf{baseados} em diagnósticos coletivos sobre a realidade daquele espaço na definição de prioridades e no compartilhamento de responsabilidades, visando a constituição de condições para o desenvolvimento integral de todos os estudantes. ( DCNEM, 2013, p. 18 ) A transformação do território deve estar articulada aos saberes locais e às possibilidades educativas, sempre com o olhar atento no suporte da aprendizagem dos estudantes, a fim de \textbf{que} o território se torne \textbf{um} lugar propício para atos educativos.
Essa intencionalidade pedagógica oportunizará a construção de práticas fundamentadas em aproximações, aceitações, diversidades de linguagens e convívios significativos.
A constituição dessas práticas pode ser efetivada com a valorização dos elementos articuladores do currículo e território, dentre eles, se destaca o reconhecimento dos saberes locais, pois auxiliam na \textbf{percepção} de práticas do dia a dia pertencentes àquele território, tais como : valores, memórias, hábitos.
Diante dessa perspectiva, percebe -se quão \textbf{imprescindíveis} são os saberes locais na construção de aprendizagens significativas e relevantes para os jovens, tendo em vista \textbf{que} possibilitam a execução de práticas e aquisição de conhecimentos a partir das vivências dos estudantes.
É mediante a esses saberes \textbf{que} podem ser criados temas geradores a \textbf{um} contexto de reflexão referente à linguagem, à matemática, às ciências humanas e às ciências da natureza.
Para \textbf{que} os saberes locais sejam reconhecidos, é necessário \textbf{que} a intencionalidade pedagógica seja articulada aos objetivos do currículo ( competências gerais, conhecimentos e habilidades interdisciplinares ou propostos pelas áreas do conhecimento ).
No currículo integrado, agentes, espaços e dinâmicas de cada território devem ser considerados como potenciais educativos capazes de proporcionar aprendizagem, haja vista \textbf{que} eles evidenciam a concretização dos saberes locais e dos objetos de reflexão das áreas do conhecimento no território vivenciado.
São compreendidos como agentes pessoas, coletivos ou instituições pertencentes ao local \textbf{que} atuam direta ou indiretamente ao realizar alterações nas suas dinâmicas a partir de intervenções.
Por meio da articulação com esses agentes há a possibilidade da remodelagem desse território, como também a formação integral dos estudantes.
Diante desse contexto, torna -se necessária a realização frequente entre os agentes e comunidade escolar por meio de diagnóstico e de escuta, os quais possibilitarão o elo com os eixos temáticos, os temas geradores ou temas estruturantes, conforme a BNCC \textbf{preconiza}.
Tal \textbf{abordagem}, direcionada às práticas de ensino e aprendizagem, confere pertinência e relevância para contextualizar os conteúdos curriculares ( unidades temáticas, objetos de conhecimento e habilidades ) no território – mediante temas, questões, saberes e práticas significativas para os estudantes.

% matched lemmas: abordagem, agregar, argumentar, autor, basear, basilar, bastante, centralidade, compromisso, condizente, configurar, conseguinte, contemporaneidade, contemporâneo, contextualização, desdobramento, detalhar, direção, disciplina, dizer, educar, ensino-aprendizagem, equidade, escopo, espírito, estímulo, explicitar, fragmentação, homem, imprescindível, inerente, integralidade, inter-relacionar, intermédio, inúmero, justo, lado, logo, mente, metodológico, norteador, percepção, permear, postar, preciso, preconizar, predominante, que, revelar, salientar, sistematização, sujeito, século, totalidade, um, uma, ver, viver
\end{textsample}
